\begin{textsample}{NEG Dim 2 – 2011\_c6 – Score -6.00 – 2011\_c6/enade\_2011\_letras\_por\_385230ec73.txt}  \label{ex:f2_neg_022}
CAPÍTULO 3 ANÁLISE TÉCNICA DA PROVA Este capítulo tem por objetivo apresentar o desempenho dos estudantes concluintes de Letras ( Bacharelado e Licenciatura ) no ENADE/2011.
Para isso, foram calculadas as estatísticas básicas da prova em seu todo, bem como as estatísticas dos componentes relacionadas à Formação Geral, ao de Conhecimento Específico da Área e das questões discursivas isoladamente.
Além disso, as análises das questões objetivas do componente de Conhecimento Específico são apresentadas separadamente para concluintes do Bacharelado e da Licenciatura, já que as questões propostas para cada grupo foram diferentes.
Nas tabelas, são apresentados o tamanho da população inscrita e de presentes, e as seguintes estatísticas das notas⁸ : média do desempenho na prova, erro padrão da média, desvio padrão, nota mínima, mediana e nota máxima.
As estatísticas apresentadas neste capítulo contemplam o total de estudantes concluintes da Área de Letras ( Bacharelado e Licenciatura ) em 2011 do Brasil e, separadamente, por Grande Região.
Foram calculadas tendo -se em vista as seguintes agregações : ( a ) as Grandes Regiões e o país como um todo ; ( b ) a Categoria Administrativa ; e ( c ) a Organização Acadêmica.
Em relação aos gráficos de distribuição de notas, o intervalo considerado foi de 10 unidades, aberto à esquerda e fechado à direita, com exceção do primeiro intervalo, [ 0 ; 10 ], fechado em ambos os extremos.
Para os gráficos de distribuição das notas das questões discursivas, foram consideradas mais duas categorias : questão em branco e nota zero.
Todos os gráficos de distribuição de notas permitem a comparação dos resultados por habilitação : Bacharelado e Licenciatura. 3.1 ESTATÍSTICAS BÁSICAS DA PROVA 3.1.1 Estatísticas Básicas Gerais A Tabela 3.1 apresenta as estatísticas básicas da prova por grande Região dos estudantes concluintes de Letras ( Bacharelado e Licenciatura ).
A população total de inscritos foi de 36.138.
Destes, 27.631 estiveram presentes, sendo 23,5 \% o índice de não comparecimento.
A Região de maior abstenção foi a Norte ( 33,0 \% ) e a de menor abstenção foi a Centro-Oeste ( 16,0 \% ). ⁸ Essas estatísticas e outras estão definidas no Capítulo 1.
A média das notas da prova como um todo ( nas seções seguintes serão analisados os componentes de Formação Geral e de Conhecimento Específico ) foi 42,6, sendo que os alunos da região Nordeste obtiveram a média mais alta ( 44,3 ) e os da região Norte obtiveram a média mais baixa ( 37,8 ).
As demais médias foram : 44,2 na região Sudeste ; 42,7 na região Sul ; e 41,0 na região Centro-Oeste.
O desvio padrão para o Brasil como um todo foi 16,2, sendo o maior desvio padrão encontrado na região Sudeste ( 16,9 ) e o menor nas regiões Norte e Centro-Oeste ( 14,7 ), indicando uma dispersão um pouco menor das notas dessas últimas regiões.
A região que obteve a maior nota máxima foi a Centro-Oeste ( 90,1 ), ao passo que a região que atingiu a menor nota máxima foi a Norte ( 85,7 ).
A mediana do Brasil como um todo foi 42,5, sendo a maior mediana obtida na região Nordeste ( 44,8 ) e a menor obtida na Norte ( 37,0 ).
A nota mínima foi zero em todas as regiões, sem exceção.
O comportamento das notas dos estudantes da Área de Letras de todo o Brasil pode ser observado no Gráfico 3.1 que apresenta a distribuição das mesmas em colunas diferentes para alunos de Bacharelado e de Licenciatura.
As distribuições são unimodais.
O intervalo modal dos concluintes de Bacharelado e de Licenciatura é o intervalo ( 40 ; 50 ].
Os coeficientes de assimetria das distribuições das notas das duas habilitações são negativos : – 0,28 para o Bacharelado e – 0,07 para a Licenciatura.
Por serem coeficientes pequenos as distribuições são quase simétricas, em especial a das notas dos alunos da Licenciatura.
A assimetria das distribuições por Grande Região varia bastante.
São negativas, concentração pouco maior do lado direito do histograma e mais espalhada do lado esquerdo, as distribuições de Licenciatura das regiões Nordeste ( – 0,25 ) e Sudeste ( – 0,21 ) e as distribuições do Bacharelado das regiões Sudeste ( – 0,25 ), Sul ( – 0,63 ) e Centro-Oeste ( – 0,71 ), esta última a mais acentuada de todas.
As demais distribuições por região e por habilitação possuem assimetria positiva, cauda mais longa do lado direito do gráfico.
Os maiores coeficientes de assimetria positiva ocorreram para as notas dos concluintes do Bacharelado das regiões Norte ( 0,80 ) e Nordeste ( 0,74 ).
Os Gráficos 3.2, 3.3 e 3.4 apresentam informações referentes à média da nota final dos Participantes, desagregando os resultados de acordo com, respectivamente, as Grandes Regiões do país, a Categoria Administrativa e a Organização Acadêmica.
Os gráficos apresentam o valor da média das notas como uma barra e os extremos do intervalo de confiança de 95 \% como linhas verticais unidas por uma linha horizontal na forma da letra H maiúscula.
Considerando -se o gráfico de notas segundo Grande Região ( Gráfico 3.2 ), observa -se que existe diferença estatisticamente significativa ao nível de 95 \% entre as duas maiores médias, obtidas nas regiões Nordeste ( 44,3 ) e Sudeste ( 44,2 ), e as médias obtidas pelas demais regiões.
Observa -se, ainda, que a menor média, obtida na região Norte ( 37,8 ) e significativamente diferentes de todas as outras.
Levando -se em conta os estudantes por Categorias Administrativas ( Gráfico 3.3 ), observa -se que existe diferença estatisticamente significante entre as médias das notas das IES Públicas ( 43,0 ) e Privadas ( 42,1 ).
Tendo como foco o Gráfico 3.4, que apresenta as notas médias das provas segundo Organização Acadêmica, constata -se que existe diferença estatisticamente significativa ao nível de 95 \% nas médias das notas dos estudantes provenientes de Centros Universitários ( 44,6 ) em relação aos concluintes de Universidades ( 42,5 ) e Faculdades ( 41,9 ). 3.1.2 Estatísticas Básicas no Componente de Formação Geral A Tabela 3.2 apresenta as estatísticas básicas em relação ao componente da prova que avalia a Formação Geral dos estudantes concluintes.
Os alunos de todo Brasil obtiveram desempenho médio de 50,5.
Quanto à variabilidade, o desvio padrão das notas dos estudantes do Brasil como um todo foi 18,5.
A maior média foi obtida nas regiões Nordeste e Sudeste ( 51,6 ), e a menor, na região Norte ( 48,0 ).
As demais médias foram : 49,9 na região Sul e 49,7 na região Centro-Oeste.
Já o maior desvio padrão foi obtido na região Sudeste ( 19,0 ) e o menor na região Centro-Oeste ( 17,7 ).
Os demais desvios padrões foram : 18,4 nas regiões Norte e Nordeste e 18,1 na região Sul.
A nota máxima ( 100,0 ) no componente de Formação Geral da prova do ENADE foi obtida por pelo menos um aluno das regiões Nordeste e Sudeste, enquanto a menor nota máxima ( 92,5 ) foi obtida nas regiões Sul e Centro-Oeste 92,5.
A mediana do Brasil como um todo foi 52,0, sendo a menor mediana encontrada na região Norte ( 48,5 ) e a maior encontrada na região Sudeste ( 54,0 ).
A nota mínima nesta parte foi zero em todas as regiões, sem exceção.
O Gráfico 3.5 propicia a avaliação do desempenho dos estudantes no componente de Formação Geral a partir do histograma da distribuição das notas correspondentes ao Bacharelado e à Licenciatura.
A distribuição de Letras – Licenciatura é unimodal, com moda no intervalo ( 50 ; 60 ], enquanto na prova como um todo ( Gráfico 3.1 ), a moda foi alcançada no intervalo imediatamente inferior, ou seja, ( 40 ; 50 ] para as duas habilitações.
Já a distribuição do bacharelado é bimodal com a moda principal coincidindo com a da Licenciatura e apresentando uma segunda moda no primeiro intervalo, [ 0 ; 10 ].
Nota -se, ainda, que no Gráfico 3.5 as notas apresentam uma maior dispersão do que no Gráfico 3.1 ( distribuição das notas da prova ), o que pode ser confirmado pela comparação dos desvios padrões : 16,2 para a nota da prova como um todo e 18,5 para o componente de Formação Geral.
Para o componente de Formação Geral, o coeficiente de assimetria da distribuição das notas dos estudantes, como na prova como um todo, também é negativo para o Bacharelado ( – 0,43 ) e para a Licenciatura ( – 0,45 ).
Apenas as distribuições do Bacharelado nas regiões Norte ( 0,43 ) e Nordeste ( 0,02 ) possuem assimetria positiva.
Em todas as demais Grandes Regiões os histogramas possuem assimetria negativa, tomando -se os alunos em conjunto ou por habilitação.
Nos Gráficos 3.6, 3.7 e 3.8 são apresentadas as informações referentes ao desempenho dos Concluintes no componente de Formação Geral, em diferentes agregações : Grande Região do país, Categoria Administrativa e Organização Acadêmica.
Observa -se pelo Gráfico 3.6 que existe diferença estatisticamente significativa entre as médias das notas no Componente de Formação Geral, segundo Grande Região do país.
Nas regiões Sudeste e Nordeste a média ( 51,6 ) é maior do que as médias das regiões Sul ( 49,9 ) e Centro-Oeste ( 49,7 ).
Da mesma forma, as médias destas últimas regiões são maiores do que a média da região Norte ( 48,0 ), a menor delas.
No Gráfico 3.7, que representa as notas médias no Componente de Formação Geral segundo Categoria Administrativa do país, observa -se que existe diferença estatisticamente significativa entre as médias.
Ao contrário do que ocorreu com as médias globais da prova, em Formação Geral os concluintes de Letras ( Bacharelado e Licenciatura considerados em conjunto ) das IES Públicas ( 49,8 ) obtiveram uma média menor do que os das IES Privadas ( 51,3 ).
Considerando -se o tipo de Organização Acadêmica, nota -se, no Gráfico 3.8, uma diferença estatisticamente significativa entre a maior média, obtida pelos concluintes de Centros Universitários ( 53,6 ) e as médias obtidas pelos alunos de Universidades ( 50,0 ) e Faculdades ( 51,3 ), sendo que a diferença entre as médias destas duas últimas categorias também é estatisticamente significativa. 3.1.3 Estatísticas Básicas do Componente de Conhecimento Específico A Tabela 3.3 apresenta as estatísticas básicas referentes ao componente de Conhecimento Específico da área de Letras ( Bacharelado e Licenciatura ).
A média do desempenho dos alunos do Brasil como um todo foi 39,9.
A maior média foi obtida nas regiões Nordeste e Sudeste ( 41,8 ), e a menor, na região Norte ( 34,4 ).
As demais médias foram : 40,2 na região Sul e 38,1 na região Centro-Oeste.
Quanto à variabilidade das notas, o desvio padrão do Brasil como um todo foi 17,6, sendo o maior desvio padrão observado na região Sul ( 18,4 ) e o menor na região Norte ( 15,8 ).
Os demais desvios foram : 17,2 na região Nordeste ; 18,3 na região Sudeste ; e 16,0 na região Centro-Oeste.
A mediana das notas dos estudantes de todo o Brasil foi 38,9.
A maior mediana ocorreu nas regiões Nordeste e Sudeste ( 41,9 ) e a menor na região Norte ( 32,3 ).
As demais medianas foram : 38,9 na região Sul e 36,9 na região Centro-Oeste.
A nota máxima do Brasil como um todo foi 92,7, sendo obtida por pelo menos um aluno da região Sudeste.
As demais notas máximas foram : 87,6 na região Norte ; 90,7 na região Nordeste ; 92,5 na região Sul ; e 91,9 na região Centro-Oeste.
A nota mínima foi zero em todas as regiões.
Assim como os Gráficos 3.1 e 3.5, o Gráfico 3.9, apresentado a seguir, proporciona uma avaliação do desempenho de concluintes dos cursos de Bacharelado e de Licenciatura em relação ao componente de Conhecimento Específico com um histograma da distribuição das notas correspondentes.
Dentre as 3 distribuições apresentadas, esta é a mais concentrada nas notas baixas.
Estas são distribuições unimodais e o grupo modal é o ( 30 ; 40 ] para ambas as opções de habilitação : Bacharelado e Licenciatura.
No entanto, para o Bacharelado nota -se outro máximo local no intervalo ( 50 ; 60 ].
No componente de Conhecimento Específico, para a distribuição das notas do curso de Bacharelado o coeficiente de assimetria é negativo e bem pequeno ( – 0,06 ) e para as de Licenciatura é positivo ( 0,13 ) e também próximo de zero.
Nota -se pelo histograma ( Gráfico 3.9 ) que para a Licenciatura a cauda é um pouco mais longa à direita e que para o Bacharelado a distribuição é praticamente simétrica.
Apenas na região Sudeste este comportamento da distribuição do Bacharelado se repete.
Nas demais regiões a assimetria é mais acentuada, sendo positiva nas regiões Norte e Nordeste e negativa nas regiões Sul e Centro-Oeste.
Para as distribuições de notas da Licenciatura os coeficientes são positivos mas praticamente nulos nas regiões Nordeste e Sudeste.
Os Gráficos 3.10, 3.11 e 3.12 apresentam uma comparação dos resultados em relação à Grande Região do país, à Categoria Administrativa e à Organização Acadêmica, agora levando em conta o desempenho de alunos no componente de Conhecimento Específico da prova.
Pelo Gráfico 3.10, observa -se que existe diferença estatisticamente significativa entre as maiores médias das notas no componente de Conhecimento Específico, das regiões Sudeste e Nordeste ( 41,8 ), em relação às demais regiões.
Nota -se, ainda, que a média da região Norte ( 34,4 ) é significativamente menor do que todas as outras.
Quanto à Categoria Administrativa ( gráfico 3.11 ), observa -se um comportamento semelhante àquele da prova como um todo, ou seja, existe diferença estatisticamente significativa entre as médias das IES Públicas e Privadas, sendo que a maior média foi obtida por alunos de IES Públicas de ensino ( 40,7 ).
Quanto ao Gráfico 3.12, observa -se, mais uma vez, que existe diferença estatisticamente significativa ao nível de 95 \% entre as notas no componente de Conhecimento Específico dos diferentes tipos de Organização Acadêmica.
A média dos Concluintes de Centros Universitários ( 41,6 ) foi maior do que de Universidades ( 40,0 ), que por sua vez foi maior do que de Faculdades ( 38,7 ). 3.2 ANÁLISE DAS QUESTÕES OBJETIVAS 3.2.1 Componente de Formação Geral A Tabela 3.4 apresenta as estatísticas básicas relativas às oito questões objetivas do componente da prova que abrange a Formação Geral dos estudantes.
A média do Brasil foi 50,3.
A menor média foi encontrada na região Norte ( 48,5 ) e a maior na região Nordeste ( 51,4 ).
As demais médias foram : 51,2 na região Sudeste ; 49,3 na região Sul ; e 49,2 na região Centro-Oeste.
O desvio padrão do Brasil foi 20,0, sendo o maior desvio padrão encontrado na região Sudeste ( 20,4 ) e o menor nas regiões Norte e Sul ( 19,7 ).
Os demais desvios foram : 20,1 na região Nordeste e 19,6 na região Centro-Oeste.
As medianas ( 50,0 ), as notas máximas ( 100,0 ) e as notas mínimas ( 0,0 ) foram iguais para todas as regiões.
A Tabela 3.5 apresenta os índices de facilidade e discriminação ( ponto bisserial ) das questões objetivas do componente de Formação Geral.
Quanto ao índice de facilidade, foram usadas as seguintes cores para diferenciar o nível de dificuldade da questão : • Azul para as questões classificadas com índice muito fácil ( índice \&gt ; =0,86 ), verde para as questões classificadas com índice fácil ( 0,61 a 0,85 ), amarelo para as questões classificadas com médio ( 0,41 a 0,60 ), vermelho para as questões classificadas com difícil ( 0,16 a 0,40 ) e roxo para as questões classificadas com muito difícil ( \&lt ; =0,15 ).
Já quanto ao índice de discriminação, foram usadas as seguintes cores para qualificar a questão : • As questões classificadas com índice fraco receberam a cor vermelho ( índice \&lt ; =0,19 ), as classificadas com médio receberam a cor amarelo ( 0,20 a 0,29 ), as classificadas com bom receberam a cor verde ( 0,30 a 0,39 ) e as classificadas com muito bom ( \&gt ; =0,40 ) receberam a cor azul.
As questões objetivas do componente de Formação Geral, segundo o índice de facilidade, foram assim avaliadas : das oito questões, nenhuma teve o índice de facilidade classificado como muito fácil.
Três questões foram tidas como fáceis, por terem índice de acertos situado na faixa entre 0,61 e 0,85 ( de 61,0 \% a 85,0 \% de acertos ).
Duas questões foram consideradas de dificuldade média, situando -se no intervalo entre 0,41 e 0,60 do índice de facilidade, ou seja, houve entre 41,0 \% e 60,0 \% de acertos.
As demais três questões foram classificadas na categoria difícil, situando -se no intervalo entre 0,16 e 0,40.
Por fim, nenhuma das questões apresentou menos de 15 \% de acertos, razão pela qual seriam classificadas como muito difíceis.
A Tabela 3.5 apresenta os índices de facilidade e discriminação ( ponto bisserial ) das questões objetivas do componente de Formação Geral.
Quanto ao índice de facilidade, foram usadas as seguintes cores para diferenciar o nível de dificuldade da questão :  Azul para as questões classificadas com índice muito fácil ( índice \&gt ; =0,86 ), verde para as questões classificadas com índice fácil ( 0,61 a 0,85 ), amarelo para as questões classificadas com médio ( 0,41 a 0,60 ), vermelho para as questões classificadas com difícil ( 0,16 a 0,40 ) e roxo para as questões classificadas com muito difícil ( \&lt ; =0,15 ).
Já quanto ao índice de discriminação, foram usadas as seguintes cores para qualificar a questão :  As questões classificadas com índice fraco receberam a cor vermelho ( índice \&lt ; =0,19 ), as classificadas com médio receberam a cor amarelo ( 0,20 a 0,29 ), as classificadas com bom receberam a cor verde ( 0,30 a 0,39 ) e as classificadas com muito bom ( \&gt ; =0,40 ) receberam a cor azul.
As questões objetivas do componente de Formação Geral, segundo o índice de facilidade, foram assim avaliadas : das oito questões, nenhuma teve o índice de facilidade classificado como muito fácil.
Três questões foram tidas como fáceis, por terem índice de acertos situado na faixa entre 0,61 e 0,85 ( de 61,0 \% a 85,0 \% de acertos ).
Duas questões foram consideradas de dificuldade média, situando -se no intervalo entre 0,41 e 0,60 do índice de facilidade, ou seja, houve entre 41,0 \% e 60,0 \% de acertos.
As demais três questões foram classificadas na categoria difícil, situando -se no intervalo entre 0,16 e 0,40.
Por fim, nenhuma das questões apresentou menos de 15 \% de acertos, razão pela qual seriam classificadas como muito difíceis.
Como já comentado, para análise das questões objetivas relativas à Formação Geral segundo o poder de discriminação, utilizou -se, o índice de discriminação ponto bisserial.
Nesta análise as questões foram assim avaliadas : cinco das oito questões apresentaram índices acima ou igual a 0,40 e, assim, foram classificadas com índice muito bom para esse grupo de estudantes.
As três outras questões tiveram bom índice de discriminação, com valor entre 0,30 e 0,39.
Nenhuma questão teve nível de discriminação médio ou fraco.
O índice de facilidade variou de 0,20 a 0,73, e o de discriminação, de 0,31 a 0,54.
As cinco questões com índices de discriminação muito bom, as de números 1, 2, 3, 5 e 6, figuraram entre diferentes níveis de dificuldade : três classificadas na categoria fácil ( questões 1, 3 e 6 ) do índice de facilidade, uma na categoria médio ( questão 5 ) e uma na categoria difícil ( questão 2 ).
Em particular, a questão 3 foi a que apresentou o maior poder discriminatório, com índice 0,54, e foi também uma das mais fáceis, com uma proporção de 0,72 acertos.
O máximo de acertos foi alcançado pela questão 6 com um índice de facilidade de 0,73.
A questão de número 4 apresentou índice de facilidade 0,53, ou seja, pouco mais da metade dos estudantes conseguiu resolvê -la, dentro do universo de participantes.
O índice de discriminação desta questão foi bom ( 0,36 ).
Já a questão 8 obteve índice de discriminação médio, 0,31, e seu índice de facilidade foi difícil ( 0,20 ).
O Gráfico 3.13, para exemplificar, analisa o comportamento da questão de número 3 de Formação Geral.
Trata -se da segunda questão mais fácil e a que obteve o maior índice de discriminação dessa parte da prova.
Neste gráfico, cada uma das cinco curvas representa o percentual de respostas em determinada alternativa da questão, em função do número de acertos dos estudantes nesta parte da prova ( Formação Geral/Múltipla Escolha ), antes de possíveis eliminações pelo critério do ponto bisserial.
Em princípio, a soma das escolhas possíveis deveria ser igual a 100 \%.
Não é este o caso, pois não aparecem no gráfico as questões deixadas em branco ou com múltiplas respostas.
Como foram oito as questões, a quantidade de acertos varia de 0 a 8.
A curva em vermelho corresponde à alternativa E, a correta para esta questão.
Assim, observa -se que entre os estudantes com menor número de acertos, nessa parte do exame, a situação mais frequente foi a escolha de uma das alternativas incorretas : a alternativa A ( em azul ).
Na medida em que o número de acertos aumenta, indicando desempenho melhor nesta parte da prova, aumenta concomitantemente a proporção de estudantes que selecionaram a alternativa correta E, atingindo 100 \% para os estudantes com 8 acertos.
Essa análise permite verificar como a questão discriminou os grupos de desempenho, justificando o alto índice obtido na questão.
Os gráficos relativos às demais questões de Formação Geral constam do Anexo I. 3.2.2 Componente de Conhecimento Específico – Letras Bacharelado A Tabela 3.6 apresenta as estatísticas básicas em relação às questões objetivas do componente de Conhecimento Específico da prova de Letras – Bacharelado, por Grande Região.
A média do Brasil deste componente foi de 42,2.
A menor média foi observada na região Norte ( 35,4 ) e a maior na região Nordeste ( 67,3 ).
O desvio padrão de todo o Brasil foi 21,9, sendo o menor desvio padrão encontrado na região Nordeste ( 9,8 ) e o maior na região Sudeste ( 27,0 ).
A mediana de todo o Brasil foi 40,9, a mesma que a da região Sudeste.
A maior mediana foi encontrada na região Nordeste ( 63,6 ) e a menor mediana ocorreu na região Norte ( 34,1 ).
A nota máxima obtida nas questões objetivas do componente de Conhecimento Específico da prova foi 95,5, de pelo menos um aluno da região Sul.
A menor nota máxima foi 68,2 obtida na região Norte.
A Tabela 3.7 apresenta os índices de facilidade e discriminação ( ponto bisserial ) das questões objetivas do componente de Conhecimento Específico para os estudantes de Letras – Bacharelado.
Para facilitar a diferenciação das questões usaram -se as mesmas cores das Tabelas 3.5 para as diferentes classificações dos índices de facilidade e de discriminação.
Dentre as questões objetivas da parte da prova relativa ao componente de Conhecimento Específico, nenhuma questão foi anulada pela Comissão.
Desse modo, a classificação quanto ao índice de facilidade foi estabelecida com base nas 27 questões ( questões 9 a 25 e 36 a 45 ).
A partir dos índices obtidos, pode -se concluir que mais da metade das questões objetivas da prova foi considerada pelo menos difícil : das 27 questões válidas, doze foram classificadas como difíceis e quatro como muito difíceis.
Não houve questão classificada como muito fácil.
Apenas uma questão foi tida como fácil, na faixa de 0,61 a 0,85 do índice de facilidade, e outras dez foram consideradas médias, entre 0,41 e 0,60.
Já quanto aos índices de discriminação das questões objetivas do componente de Conhecimento Específico da prova, tem -se como resultado a seguinte classificação : dezessete das 27 questões foram consideradas como muito boas, enquanto quatro delas tiveram índice de discriminação bom.
Assim, para a maioria das questões – 21 em 27 – os índices de discriminação foram bons ou muito bons.
Dentre as demais, uma delas foi classificada como média e as outras cinco como fracas, sendo seis, por conseguinte, a quantidade de questões nos dois patamares mais baixos de discriminação.
Constata -se, assim, que a parte objetiva da prova – no que se refere ao componente de Conhecimento Específico do Bacharelado – possuía capacidade média de discriminar entre aqueles que dominam ou não o conteúdo.
Dentre as dezessete questões que alcançaram os maiores índices de discriminação, classificadas com índice de discriminação muito bom, situando -se no intervalo de 0,40 a 0,50 do índice, apenas uma delas ( questão 12 ) foi classificada na categoria fácil, quanto ao índice de facilidade.
Dentre as demais questões com discriminação muito boa, oito foram classificadas como de dificuldade média ( questões 11, 13, 22, 25, 36, 38, 43 e 45 ) e outras oito foram classificadas com difíceis ( questões 17, 18, 19, 20, 21, 23, 39 e 41 ).
A questão de número 40 foi a mais difícil dentre as 27 questões específicas válidas, com baixo índice de facilidade, apenas 9,0 \% de acertos.
Essa questão apresentou poder discriminatório igualmente baixo, 0,06, o que comprova ter sido esta a mais difícil para os estudantes.
Destaca -se, também, a questão 44, com índice de facilidade 0,12, o que, em termos percentuais, corresponde a 12,0 \% de estudantes que responderam acertadamente.
Foi também muito fraco seu índice de discriminação, 0,07.
Tais questões foram, portanto, pelo critério ponto bisserial consideradas inadequadas.
Por isso, as questões 40 e 44 foram eliminadas do cômputo da nota final.
Além destas duas, as demais questões com índice fraco de discriminação, questões 14, 16 e 42 também não foram utilizadas no cômputo final das notas, num total de 5 questões eliminadas.
A questão 42 chegou a apresentar um índice de discriminação negativo, -0,05, ainda que bem perto de zero.
A título de exemplo das análises do comportamento das questões objetivas, o Gráfico 3.14 analisa a questão 25 do componente de Conhecimento Específico.
Esta foi a questão com maior poder de discriminação ( 0,60 ) e considerada de dificuldade média, com 57,0 \% de estudantes que assinalaram acertadamente a opção B, correspondente ao gabarito.
Neste gráfico, cada uma das cinco curvas representa o percentual de respostas em determinada alternativa da questão 25, em função do número de acerto dos estudantes nesta parte da prova, antes de possíveis eliminações de questões pelo critério do ponto bisserial.
A alternativa correta B, representada no gráfico pela curva em verde, foi escolhida em maiores proporções pelos alunos com desempenho melhor nesta parte da prova.
Já as alternativas incorretas, também denominadas distratores, foram selecionadas principalmente por aqueles com notas mais baixas.
Neste caso também a soma não é sempre 100 \% por causa das questões não respondidas ou com mais de uma opção marcada.
Aqueles com nota zero, na sua quase totalidade deixaram esta questão em branco ou marcaram mais de uma alternativa, comportamento considerado inválido.
A proporção de alunos que selecionou a resposta correta B aumenta gradativamente, chegando a atingir 100 \% para alunos com mais de 20 acertos, enquanto a proporção dos que escolheram alternativas incorretas decai a partir de 4 ou 5 acertos como função do número de acertos nesta parte da prova.
Os gráficos relativos às demais questões do componente de Conhecimento Específico constam do Anexo I. 3.2.3 Componente de Conhecimento Específico – Letras Licenciatura A Tabela 3.8 apresenta as estatísticas básicas em relação às questões objetivas do componente de Conhecimento Específico da prova de Letras - Licenciatura, por Grande Região.
A média do Brasil deste componente foi de 42,4.
A menor média foi observada na região Norte ( 37,2 ) e a maior na região Nordeste ( 44,6 ).
O desvio padrão de todo o Brasil foi 18,6, sendo o maior desvio padrão encontrado na região Sul ( 19,1 ) e o menor na região Norte ( 17,0 ).
A mediana de todo o Brasil foi 41,7, a mesma encontrada na região Sul.
A maior mediana, 45,8, ocorreu nas regiões Nordeste e Sudeste e a menor, 37,5, nas regiões Norte e Centro-Oeste.
A nota máxima da prova foi 100,0, obtida por pelo menos um concluinte de Licenciatura da região Nordeste, nas questões objetivas do componente de Conhecimento Específico.
As notas máximas nas demais regiões foram : 91,7 na região Norte ; 95,8 nas regiões Sudeste e Sul ; e 91,7 na região Centro-Oeste.
A Tabela 3.9 apresenta os índices de facilidade e discriminação ( ponto bisserial ) das questões objetivas do componente de Conhecimento Específico da Área de Letras - Licenciatura.
Para facilitar a diferenciação das questões usaram -se as mesmas cores das Tabelas 3.5 e 3.7 para as diferentes classificações dos índices de facilidade e de discriminação.
Dentre as questões objetivas da parte da prova relativa ao componente de Conhecimento Específico ( questões 9 a 35 ), nenhuma foi anulada pela Comissão.
Desse modo, a classificação quanto ao índice de facilidade foi estabelecida com base em todas as 27 questões.
A partir dos índices obtidos, pode -se concluir que mais da metade das questões objetivas da prova foi considerada pelo menos difícil : das 27 questões, quinze foram classificadas como difíceis e duas como muito difíceis.
Não houve questão classificada como muito fácil, ao passo que três foram tidas como fáceis, na faixa de 0,61 a 0,85 do índice de facilidade, e outras sete consideradas médias, entre 0,41 e 0,60.
Já quanto aos índices de discriminação das questões objetivas do componente de Conhecimento Específico da prova, tem -se como resultado a seguinte classificação : dez das 27 questões foram consideradas como muito boas, enquanto outras dez tiveram índice de discriminação bom.
Assim, para a maioria das questões – 20 em 27 – os índices de discriminação foram bons ou muito bons.
Dentre as demais, quatro delas foram classificadas como médias e três como fracas, sendo sete, por conseguinte, a quantidade de questões nos dois patamares mais baixos de discriminação.
Constata -se, assim, que as questões objetivas da prova de Licenciatura – no que se refere ao componente de Conhecimento Específico – possuíam boa capacidade de discriminar entre aqueles que dominam ou não o conteúdo.
Dentre as questões que alcançaram os maiores índices de discriminação, classificadas com índice muito bom, situando -se no intervalo de 0,42 a 0,50 do índice, duas delas ( questões 12 e 27 ) foram classificadas na categoria fácil, cinco ( questões 22, 25, 32, 33 e 34 ) como médias e outras três ( questões 11, 13 e 18 ) como difíceis, quanto ao índice de facilidade.
A questão de número 16 foi uma das mais difíceis dentre as 27 questões específicas, com baixo índice de facilidade, apenas 15,0 \% de acertos.
Essa questão apresentou poder discriminatório igualmente baixo, 0,07, o que comprova ter sido esta a mais difícil para os estudantes.
Destaca -se, também, a questão 26, com índice de facilidade 0,24, o que, em termos percentuais, corresponde a 24,0 \% de estudantes que responderam acertadamente, obtendo, ainda, 0,09 de índice de discriminação.
As questões 16 e 26 foram eliminadas do cômputo da nota final pelo critério do índice de discriminação ponto bisserial.
Além destas duas, a questão 14, por ter índice fraco de discriminação, também não foi utilizada no cômputo da nota final.
A título de exemplo das análises do comportamento das questões objetivas, o Gráfico 3.15 foca a questão 25 do componente de Conhecimento Específico.
Esta foi uma das questões comuns para o Bacharelado e a Licenciatura e, da mesma forma que foi observado para os alunos de Bacharelado, esta questão foi a de maior capacidade de discriminação com índice de discriminação 0,58 para os alunos da Licenciatura, apresentando índice de facilidade 0,56, ou seja, 56,0 \% dos estudantes assinalaram acertadamente a opção B, correspondente ao gabarito.
Neste gráfico, cada uma das cinco curvas representa o percentual de respostas em determinada alternativa da questão 25, em função do número de acertos dos estudantes nesta parte da prova, antes de possíveis eliminações de questões pelo critério do ponto bisserial.
A alternativa correta B, representada no gráfico pela curva em verde, foi escolhida em maiores proporções pelos alunos com desempenho melhor nesta parte da prova.
Já as alternativas incorretas, também denominadas distratores, foram selecionadas principalmente por aqueles com notas mais baixas.
Aqueles com nota zero, na sua quase totalidade deixaram esta questão em branco ou marcaram mais de uma alternativa, comportamento considerado inválido.
A proporção de alunos que selecionou a resposta correta B aumenta gradativamente, chegando a atingir 100 \% para 22 acertos ou mais, enquanto a proporção dos que escolheram alternativas incorretas decai, a partir de 3 ou 4 acertos, como função do número de acertos nesta parte da prova.
Os gráficos relativos às demais questões do componente de Conhecimento Específico constam do Anexo I. 3.3 ANÁLISE DAS QUESTÕES DISCURSIVAS 3.3.1 Componente de Formação Geral As análises dos resultados de desempenho dos estudantes de Letras ( Bacharelado e Licenciatura ) nas duas questões discursivas relativas à Formação Geral encontram -se na Tabela 3.10 e no Gráfico 3.16.
Na tabela 3.10 observa -se que a nota média nesse conjunto de questões foi bastante próxima da obtida nas objetivas.
Os estudantes de todo o Brasil obtiveram, em Formação Geral, média 50,3 nas questões objetivas e 50,9 nas questões discursivas.
No entanto, pode -se notar um aumento do desvio padrão de 20,0 nas questões objetivas do componente de Formação Geral dos alunos de todo o Brasil, para 27,8 nas questões discursivas do mesmo componente.
A maior média foi obtida na região Sudeste ( 52,2 ) e a menor na região Nordeste ( 47,2 ).
A mediana de todo o Brasil, neste componente, foi 55,0, a mesma das regiões Nordeste e Sudeste.
Já nas regiões Sul e Centro-Oeste a mediana foi maior, 52,5, e na região Norte foi menor, 50,0.
A nota máxima ( 100,0 ) e a nota mínima ( 0,0 ) foram as mesmas em todas as regiões do Brasil.
O Gráfico 3.16 representa a distribuição das notas nas questões discursivas no componente de Formação Geral segundo a opção de Bacharelado/Licenciatura.
A moda da distribuição das notas dos estudantes de Bacharelado ocorre no intervalo [ 0 ; 10 ], intervalo que inclui, além das notas zero, a frequência de alunos que deixaram este tipo de questão em branco.
Para a habilitação em Licenciatura a moda é o intervalo ( 40 ; 50 ] e observam -se outros máximos locais nos intervalos [ 0 ; 10 ] e ( 70 ; 80 ].
Para o Bacharelado, os máximo locais encontram -se em ( 40 ; 50 ] e ( 70 ; 80 ].
Os coeficientes de assimetria das duas distribuições são negativos para o Brasil como um todo : – 0,13 para o Bacharelado e – 0,43 para a Licenciatura.
As distribuições de ambos os cursos em todas as regiões tem o mesmo comportamento, são assimétricas à esquerda, possuem cauda mais acentuada no lado esquerdo do gráfico.
Exceção disso ocorre apenas para o curso de Bacharelado da região Nordeste cuja distribuição possui forte assimetria positiva ( coeficiente 1,22 ).
Na sequência, os resultados verificados para cada uma das questões discursivas de Formação Geral serão apresentados, estabelecendo -se relações com os conteúdos abordados em cada uma delas.
Os comentários da Banca de docentes corretores a respeito do observado na correção das respostas dos estudantes, suas impressões e conclusões serão apresentados junto à análise de cada questão.
Cumpre esclarecer que, tendo em vista que as questões discursivas de Formação Geral são padronizadas, ou seja, constam de todas as provas, os comentários da Banca são os mesmos para todas as carreiras acadêmicas, sendo direcionados a todos os estudantes que participaram do ENADE/2011.
A seguir, serão analisados os desempenhos da Área de Letras ( Bacharelado e Licenciatura ) nas duas questões discursivas de Formação Geral do ENADE/2011, comparando os resultados obtidos com comentários para cada questão. 3.3.1.1 Análise da Questão Discursiva 1 do Componente de Formação Geral Os dados de Letras ( Bacharelado e Licenciatura ), obtidos a partir das respostas à questão 1, encontram -se na Tabela 3.11 e no Gráfico 3.17.
Nessa questão – de melhor desempenho dentre as duas de Formação Geral – os alunos de todo Brasil tiveram média, 55,4.
A maior média para a questão 1 foi obtida na região Sudeste ( 57,9 ), e a menor, na região Norte ( 51,6 ).
Quanto à variabilidade das notas, o desvio padrão de todo o Brasil foi 35,8.
O menor desvio padrão foi obtido na região Centro-Oeste ( 35,3 ) e o maior desvio padrão foi obtido na região Norte ( 36,2 ).
As medianas das regiões Nordeste, Sul e Centro-Oeste, bem como do Brasil como um todo, foram iguais ( 60,0 ).
Na região Sudeste a mediana foi 65,0 e na região Norte foi 50,0.
As notas máximas e mínimas da questão discursiva 1 foram as mesmas para todas as regiões do Brasil, respectivamente, 0,0 e 100,0.
O Gráfico 3.17 mostra a distribuição das notas na questão discursiva 1 do componente de Formação Geral para alunos do Bacharelado e da Licenciatura.
Observa -se que a maior frequência corresponde aos alunos que obtiveram notas no intervalo ( 90 ; 100 ], tanto para os de Bacharelado quanto para os de Licenciatura.
No entanto, é significativa a porcentagem de alunos do Bacharelado que não resolveu esta questão.
A menos da distribuição de notas dos concluintes de Bacharelado da região Centro-Oeste ( 0,17 ), o coeficiente de assimetria é negativo para as duas habilitações em todas as regiões.
Na maioria dos casos os coeficientes têm valor absoluto pequeno, variando de – 0,03 ( Licenciatura da região Norte ) a – 0,35 ( Licenciatura da região Sudeste ).
Apenas para as notas dos alunos de Bacharelado da região Nordeste a assimetria negativa é mais acentuada ( com um coeficiente de – 1,97 ). 3.3.1.2 Comentários sobre a correção das respostas à Questão Discursiva 1 De maneira geral, a aplicação da chave de correção da questão 1 de Formação Geral não apresentou qualquer dificuldade digna de menção.
Para isso contribuíram, indubitavelmente, os ajustes feitos após a correção da amostra e a reunião entre todos os membros da Banca de docentes corretores.
As poucas dúvidas, todas pontuais, apresentadas pelos corretores, foram acompanhadas e respondidas pela coordenação e subcoordenação da correção das questões de Formação Geral, por meio da ferramenta de Gerenciamento de Dúvidas do Sistema de Correção On-line.
Não houve registro de qualquer ocorrência que pusesse em xeque o padrão de resposta ou a efetividade e a adequação da chave de correção.
Explica -se : trata -se de questão com comando claro, direto e objetivo ( solicitava -se, basicamente, três vantagens justificadas de cursos a distância ), cujas respostas foram corrigidas por meio da aplicação de um chave de correção testada e aprovada previamente.
Havia absoluta clareza quanto aos critérios de avaliação da correspondência entre as respostas dos estudantes e as possibilidades de vantagens de cursos a distância admitidas como corretas no padrão de resposta oficial, além de gradações explícitas ( e fáceis de aplicar ) dos diferentes níveis de pontuação previstos.
Felizmente, portanto, não há reparo a registrar em relação à facilidade de aplicação do padrão de resposta e da chave de correção, e nem em relação à atribuição dos diferentes níveis de pontuação previstos.
Todas as dificuldades que poderiam ter obstado a correta aplicação do padrão de resposta oficial e da respectiva chave de correção foram evitadas por meio dos ajustes feitos após a correção da amostra e fartamente debatidos com toda a Banca.
Digno de nota é que quantidade tão significativa de profissionais envolvidos na mesma tarefa – tanto para a questão 1 quanto para a 2 – tenha apresentado tão poucas dificuldades na execução da correção, em termos proporcionais.
Em suma, a correção da questão 1 da prova de Formação Geral do ENADE 2011 foi exemplarmente bem planejada, servindo -se de padrão de resposta muito bem adequado à questão proposta.
Quanto ao tema desta questão, em particular, a Banca verificou que uma parcela significativa de estudantes evocou experiências bastante concretas e próximas de sua realidade.
Houve várias respostas que indicavam uma vivência pessoal de ensino superior na modalidade Educação a Distância ( EaD ), evidenciando o tom de depoentes nos textos apresentados.
Os estudantes, em sua grande maioria, utilizaram parte considerável do espaço de 15 linhas disponíveis para a resposta – e outra parcela menos significativa dos que não o fizeram demonstraram notável capacidade de atender ao comando da questão de maneira objetiva, curta e, via de regra, correta.
Registre -se, ainda, que foram relativamente poucos os casos de respostas que tenham passado ao largo do tema em pauta na questão.
As capacidades de leitura, de compreensão do comando proposto e de expressão escrita que os estudantes avaliados na edição 2011 do ENADE foram satisfatórias.
Não obstante, seria leviano perder de vista que a qualidade dos textos redigidos em resposta às questões discursivas do Exame ainda está muito aquém do que se espera de concluintes de cursos de ensino superior de todas as regiões do país.
Quanto ao conteúdo das respostas, a Banca constatou boa capacidade, por parte da maioria dos estudantes, de compreensão do tema e do comando da questão.
Foram relativamente poucos os casos de respostas que deixaram de enumerar vantagens da modalidade EaD, e proporcionalmente escassos os estudantes que citaram vantagens não previstas no padrão de resposta.
Foi frequente, a tentativa direta de atender ao comando da questão.
Os erros mais comuns, em relação ao padrão de respostas e à grade de correção, foram fruto do desdobramento em vários “ itens ” daquilo que, de acordo com o padrão de respostas oficial, representava uma única vantagem.
Destacaram -se, neste caso, as respostas que apontavam a flexibilidade de horário e/ou local como duas vantagens distintas daquela modalidade de ensino.
Quanto aos diferentes níveis de pontuação previstos, a maior causa de baixas pontuações foi a ausência de justificativas, e mesmo de argumentação, para uma ou mais das vantagens enumeradas.
Isso demonstra não apenas a objetividade das respostas, coerente com a objetividade do comando da questão ( “ enumere três vantagens de um curso a distância ” ), mas também certa dificuldade de formulação plena de um texto, ou ao menos de parágrafos, em formato dissertativo – mesmo diante de uma média de 5 linhas disponíveis para cada vantagem a enumerar.
Os acertos mais comuns, ou seja, os “ itens ” do padrão de respostas mais frequentemente mencionados foram : ( 1 ) a flexibilidade de horário e/ou local ; ( 2 ) a capilaridade do ensino a distância ; ( 3 ) a democratização do acesso à educação de qualidade ; e ( 4 ) os custos menores que os de cursos presenciais.
Dentre as vantagens previstas no padrão de respostas que foram menos citadas, destacam -se a inclusão de pessoas com comprometimento motor, a qualificação de professores e a troca de experiências entre os participantes.
Foram muito comuns, no entanto, as menções à supostamente maior facilidade de acesso a professores ou/e tutores em cursos superiores a distância. 3.3.1.3 Análise da Questão Discursiva 2 do Componente de Formação Geral A Tabela 3.12 mostra que o desempenho dos estudantes na questão 2 ( média 46,2 ) foi inferior ao obtido na questão de número 1 ( média 55,4 ).
A região Norte foi aquela onde a média, nessa questão, foi menor ( 42,7 ), e a de maior média foi a região Nordeste ( 48,2 ).
Quanto à variabilidade das notas, o desvio padrão de todo o Brasil foi 30,9, inferior ao obtido na questão de número 1 ( 35,8 ).
O maior desvio nessa questão foi obtido na região Sudeste ( 31,3 ), enquanto o menor foi obtido na região Sul ( 30,3 ).
As medianas foram 50,0 para todas as regiões, exceto na região Nordeste, onde esta estatística foi maior, 60,0.
A nota máxima ( 100,0 ) foi obtida em todas as regiões do Brasil, bem como a nota mínima 0,0.
O Gráfico 3.18 mostra a distribuição das notas na questão discursiva 2 do componente de Formação Geral para alunos do Bacharelado e da Licenciatura.
Observa -se que a maior frequência corresponde aos alunos que deixaram a questão em branco, tanto para aqueles de Bacharelado quanto para os de Licenciatura, com frequência mais elevada para o Bacharelado.
Nota -se, ainda, que se considerássemos apenas os alunos que resolveram esta questão, a moda seria o intervalo ( 50 ; 60 ] para as duas habilitações.
O coeficiente de assimetria da distribuição de notas do Bacharelado é positivo ( 0,17 ) e da Licenciatura negativo ( – 0,33 ).
Em todas as regiões este comportamento se repete : assimetria positiva para a distribuição do Bacharelado e negativa para a da Licenciatura.
Apenas na região Norte a assimetria é negativa para as duas habilitações. 3.3.1.4 Comentários sobre a correção das respostas à Questão Discursiva 2 Diferentemente da questão 1, a questão 2, cujo tema aborda políticas públicas para a erradicação do analfabetismo e para o aumento da empregabilidade, permitiu mais amplitude na elaboração das respostas dos estudantes, uma vez que solicita a apresentação de uma proposta.
À falta de objetividade técnica na elaboração da questão – sobretudo no seu comando – correspondeu uma miríade de possibilidades interpretativas.
Registre -se que não raras foram as respostas que se utilizaram de lugares comuns e exortações religiosas/humanitárias/cívicas, contudo, na grande maioria dos casos, essas exortações foram usadas como um encerramento do texto e não comprometeram a resposta.
Muito frequente também foi a confusão entre política pública e política partidária, bem como alguns poucos, confundiam programas educacionais com programas educativos nas TVs.
Assim, destacaram -se como propostas/programas mais recorrentes :  prosseguimento das políticas já existentes, mas com o aumento dos investimentos, normalmente sugerindo a ampliação da bolsa família, da Educação de Jovens e Adultos ( EJA ) e da EaD ;  valorização do magistério, construção de escolas e melhoria das atuais ;  parceria do governo com empresas para manutenção de salas de aula visando a alfabetização de seus funcionários, oferecimento de estágios e redução da carga horária em troca de \textbf{isenção} de impostos ;  parcerias com igrejas e ONGs para criação de espaços de alfabetização ;  escolas itinerantes e alfabetizadores em domicílio, principalmente para pessoas com dificuldade de locomoção, como os idosos e deficientes físicos ;  erradicação do trabalho infantil ;  vinculação da bolsa família não apenas à frequência, mas também e, principalmente, ao resultado obtido pelo aluno na escola ;  revisão das políticas atuais, sendo a mais frequente a extinção da aprovação automática ;  críticas consistentes em relação a modelos didáticos considerados inadequados e desestimulantes para a educação não só de adultos, mas de pessoas de todas as idades.  Destacam -se, ainda, outras sugestões apresentadas :  \textbf{Creches} nas escolas onde os pais estão sendo alfabetizados.  Diminuição de duas horas na jornada de trabalho em empresas para funcionários não alfabetizados, para que possam frequentar a escola.  Campanhas educativas vinculadas aos meios de comunicação.  Aumento do número de escolas noturnas.  Formação específica para professores alfabetizadores.
Quanto à relação entre o analfabetismo e a empregabilidade, deve -se sublinhar que nem todos os estudantes estabeleceram claramente o vínculo entre essas duas situações sociais.
Alguns falaram separadamente de uma e de outra.
Mas a maioria fez referência à necessidade de estudo para “ conseguir um bom emprego com um bom salário ”.
Alguns, em menor número, estabeleceram de forma bastante interessante a questão histórica para a situação do Nordeste ; e, também, a relação entre a pessoa analfabeta/com pouca instrução e a desempregada/empregada em serviços mais pesados e pior remunerados, que não teriam condições de educar seus filhos que, por sua vez, também não teriam melhores oportunidades no mercado de trabalho, identificando a formação de um círculo \textbf{vicioso} e a necessidade de sua interrupção.
Muitos afirmaram que o analfabetismo não é o único responsável pelo desemprego e sim a má distribuição de renda.
Grande também foi o número de alunos que criticaram o resultado da pesquisa, afirmando que a mesma é enganosa, uma vez que considera alfabetizada a pessoa que “ desenha seu nome ”, sendo comum associarem a atual política de alfabetização com ganhos eleitorais.
Foi comum a resposta incompleta, em que o estudante fez apenas a análise das desigualdades/crítica do quadro apresentado, ou só apresentou proposta.
Alguns se limitaram a responsabilizar o governo referindo -se de forma bastante genérica à questão da “ educação ” e “ profissional ”.
Entre os equívocos que mais se repetiram, destacou -se a simples análise dos dados apresentados na tabela que consta do enunciado.
Alguns textos, inclusive, estavam corretos, sem que, no entanto, fosse respondida a pergunta.
Da mesma forma, quando os estudantes partiram para a segunda parte da questão, fizeram referência a vários programas já existentes ou simplesmente disseram que é muito importante que existam projetos para a educação.
Outro equívoco recorrente foi a análise da educação no Brasil como um todo.
Essa análise, apesar de correta, não se referia especificamente à questão do analfabetismo, o que, por conseguinte, levava a sugestões que não eram direcionadas à erradicação do analfabetismo, e sim à melhora da educação no Brasil.
Sendo assim, foram apresentadas propostas como o aumento do número de faculdades, o reforço de alunos do Ensino Médio, a criação de escolas técnicas, etc.
Por outro lado, houve um grande número de redações bem escritas e precisas, no que se refere ao que foi exigido pela questão.
Foram análises equilibradas e sensatas, correlacionando a problemática do analfabetismo com o desemprego e a apresentação de sugestões bastante consistentes. 3.3.2 Componente de Conhecimento Específico Na parte da prova relativa às questões discursivas no componente de Conhecimento Específico ( Tabela 3.13 ), observa -se que a média foi mais baixa do que para as questões discursivas do componente de Formação Geral.
Enquanto no componente de Formação Geral a média para estudantes de Letras de todo o Brasil foi 50,9, na parte de Conhecimento Específico a média foi 26,0.
A maior média deste componente foi obtida pelos estudantes da região Sudeste ( 29,5 ), e a menor, pelos da região Norte ( 18,2 ).
Quanto à variabilidade das notas, o desvio padrão de todo o Brasil foi 23,6.
O maior desvio padrão foi encontrado na região Sudeste ( 24,9 ), e o menor, na região Norte ( 20,3 ).
A maior nota máxima, 100,0, foi obtida por pelo menos um aluno da região Nordeste, enquanto a menor nota máxima foi encontrada na região Norte ( 90,0 ).
Além disso, a nota mínima ( 0,0 ) foi obtida por alunos de todas as regiões do Brasil sem exceção.
A mediana do Brasil como um todo foi 23,3.
A maior mediana foi encontrada nas regiões Sudeste ( 28,3 ) e a menor ocorreu na região Norte ( 13,3 ).
O Gráfico 3.19 representa a distribuição das notas nas questões discursivas no componente de Conhecimento Específico segundo a opção de Bacharelado/Licenciatura.
As modas destas distribuições ocorrem no primeiro intervalo, [ 0 ; 10 ].
O percentual de alunos de Licenciatura com notas no intervalo [ 0 ; 10 ] é menor do que o de Bacharelado, ( cerca de 40 \% para o Bacharelado e cerca de 35 \% para a Licenciatura ) enquanto esta proporção se inverte para os intervalos até 50 pontos e voltando a inverter nos intervalos seguintes.
Nota -se que ambas as distribuições são decrescentes e que no intervalo ( 90 ; 100 ] ocorrem pouquíssimas notas.
Este tipo de comportamento das distribuições, assimetria positiva, é confirmado por coeficientes positivos : 0,53 para o Bacharelado e 0,49 para a Licenciatura.
O único caso de distribuição assimétrica à esquerda ocorre para as notas dos alunos de Bacharelado da região Centro-Oeste.
Em todas as outras agregações por região e habilitação os coeficientes de assimetria são positivos.
A análise de cada uma destas questões será feita a seguir. 3.3.2.1 Análise da Questão Discursiva 3 do Componente de Conhecimento Específico Na questão 3, cujos resultados aferidos encontram -se descritos na Tabela 3.14, a média dos estudantes de todo o Brasil foi 34,8.
A menor média nessa questão foi obtida pelos alunos da região Norte ( 25,8 ), enquanto a maior média foi obtida na região Sudeste ( 37,9 ).
Quanto à variabilidade das notas, o desvio padrão de todo o Brasil foi 29,8.
O maior desvio padrão foi obtido na região Sudeste ( 30,5 ), enquanto o menor foi obtido na região Norte ( 27,8 ).
A nota máxima, 100,0 pontos, foi alcançada por pelo menos um aluno de todas as regiões, da mesma forma que a nota mínima zero.
A mediana do Brasil como um todo foi 40,0, e foi a mesma nas regiões Sul e Centro-Oeste.
Nas demais regiões as medianas foram : 45,0 nas regiões Sudeste e Sul e 20,0 na região Norte.
O Gráfico 3.20 mostra a distribuição das notas na questão discursiva 3, do componente de Conhecimento Específico da área de Letras para alunos do Bacharelado e da Licenciatura.
Essa distribuição tem moda nas questões em branco, com percentual maior para os alunos do Bacharelado.
Nota -se ainda um máximo local no intervalo ( 40 ; 50 ] para ambas as habilitações.
O coeficiente de assimetria é pequeno para os dois cursos : 0,16 para o Bacharelado e 0,05 para a Licenciatura.
A assimetria é positiva para a distribuição das duas habilitações nas regiões Sul e Norte ; é negativa para os dois cursos na região Nordeste ; na região Centro-Oeste é positiva apenas para a Licenciatura ; e na região Sudeste é positiva apenas para o Bacharelado. 3.3.2.2 Comentários sobre a correção das respostas à Questão Discursiva 3 A questão 3 apresenta dois textos motivadores : um fragmento de ensaio de Umberto Eco, segundo o qual a interpretação do texto literário não é totalmente livre, mas deve se ater à intenção do texto ; e uma canção de Gilberto Gil, na qual se apresenta a tese de que as ideias do autor do texto literário não podem ser tão precisas, devido à existência de ambiguidade, de linguagem metafórica ( o termo “ metáfora ” é, inclusive, o título da música ).
Em seguida, solicita -se que o estudante elabore um texto dissertativo a respeito das “ especificidades da linguagem literária ”.
Tal solicitação, um tanto abrangente, se apresenta em \textbf{negrito}, centralizada na página.
Na linha de baixo, sem destaque e sem centralização, determina -se que se deve traçar um paralelo entre dois pontos específicos : a “ intenção do texto ” e a “ liberdade de interpretação ”.
A questão foi formulada de modo a dar maior destaque, portanto, à exposição sobre as especificidades da linguagem literária do que à discussão sobre liberdade de expressão X intenção do texto.
O comando solicita, explicitamente, que seja elaborado um texto dissertativo sobre o tema da questão, mas não recomenda que o concluinte utilize o padrão formal culto da língua portuguesa.
A temática e o grau de dificuldade são compatíveis com o nível de conhecimento esperado de um graduando em Letras.
Como o enunciado da questão deu mais ênfase à abordagem das especificidades da linguagem literária, uma parte dos estudantes dissertou sobre esse aspecto, sem elaborar uma reflexão a respeito da relação ‘ intenção do texto ’ vs. ‘ liberdade de interpretação ’.
Por isso, houve penalização da resposta que apenas elencou características da linguagem literária sem explorar as duas possibilidades de abordagem do texto literário.
Como o comando solicitava a elaboração de um texto dissertativo, não houve penalização da resposta que realizou apenas uma exposição da dialética ‘ intenção do texto ’ vs. ‘ liberdade de interpretação ’, sem apoiar, argumentativamente, uma ou outra visão.
No que diz respeito ao conteúdo, foram poucos os graduandos que conseguiram relacionar satisfatoriamente a dialética exigida.
Houve um grande número de estudantes que, descartando o conteúdo “ intenção do texto ”, desenvolveram, prioritariamente, apenas o tópico da metáfora e/ou da liberdade de interpretação ; outros abordaram o tema solicitado de forma pouco aprofundada, sem levantar argumentos que sustentassem sua opinião.
A julgar pela maioria das respostas, um dos problemas apresentados foi a ausência de argumentação, revelando que os Cursos de Letras não estão dando a devida atenção ao desenvolvimento de estratégias textuais de caráter argumentativo para que possam refletir sobre um determinado assunto a partir do contraponto de ideias.
Nessas respostas, não se verificou preocupação em discutir o problema, mas sim em fazer afirmações destituídas de argumentos que as comprovassem.
Na verdade, o concluinte apresentou opiniões prontas.
Essa ausência de argumentação revela o desconhecimento de questões teóricas básicas sobre a linguagem literária.
Houve, ainda, casos em que os alunos discorreram sobre características dos textos em geral, e não especificamente do texto literário.
Um importante aspecto a destacar é o mau desempenho da grande maioria dos estudantes em relação à estrutura formal do texto produzido ( mecanismos de coesão e utilização do padrão formal culto ), o que é extremamente preocupante ao se levar em conta que são formandos de Letras.
Foram muito comuns os casos de erros de concordância ( mesmo com o verbo posposto ao sujeito ), pontuação, acentuação, regência, seleção vocabular e estruturação frasal, com uso inadequado ou ausência de conectivos entre parágrafos e entre frases.
Os textos careceram de um mínimo de textualidade e de domínio do registro padrão da língua.
Na verdade, observaram -se relações linguísticas quase agramaticais, como as estabelecidas pela sequência de gerúndios sem o apoio de um ponto de partida para a organização das informações gramaticais e semânticas.
Conclui -se, portanto, que pouquíssimos conseguiram de fato desenvolver os tópicos solicitados e relacioná -los produtivamente.
Houve um grande número de respostas que caíram no lugar-comum da liberdade de interpretação ; muitos casos de interpretação insatisfatória de um ou dos dois tópicos, havendo várias respostas de pouquíssimas linhas.
De fato, poucos conseguiram desenvolver esses tópicos a contento, e menos ainda relacioná -los.
Nas respostas que apresentaram bom desempenho, tanto a forma quanto o conteúdo foram satisfatoriamente trabalhados : no que diz respeito ao conteúdo, foram abordados os três tópicos fundamentais da questão, houve desenvolvimento da ideia de “ intenção do texto ”, das “ múltiplas possibilidades de leituras ” e da relação dialética entre elas ; quanto à forma do texto, não foram observadas incorreções gramaticais nem ortográficas e houve paragrafação pertinente e adequado uso de conectivos.
Nas respostas em que o desempenho foi mediano, houve pequenas incorreções na forma ou no conteúdo.
No que diz respeito ao conteúdo, os graduandos abordaram os três tópicos fundamentais com alguma superficialidade ou omitiram algum deles, houve pouco desenvolvimento da contraposição entre “ intenção do texto ” e “ múltiplas possibilidades de leituras ”.
No que diz respeito à forma do texto, observou -se incorreções gramaticais de pouca importância ( pontuação, concordância, acentuação e regência ) ou problemas na paragrafação e no uso de conectivos.
Respostas de fraco desempenho foram as que tiveram incorreções na forma ou no conteúdo.
A respeito do conteúdo, os alunos abordaram com bastante superficialidade o contraste entre “ intenção do texto ” e “ múltiplas possibilidades de leituras ”, não houve detalhamento das especificidades da linguagem literária, principalmente a questão da ambigüidade.
No que diz respeito à forma do texto, observou -se muitas incorreções gramaticais ( pontuação, concordância, acentuação e regência ) e texto truncado por problemas na paragrafação e no uso de conectivos.
Por fim, as respostas que apresentaram desempenho insuficiente foram as que tiveram muitas incorreções na forma ou no conteúdo : no que diz respeito ao conteúdo, a abordagem do contraste entre “ intenção do texto ” e “ múltiplas possibilidades de leituras ”, foi equivocada, superficial e/ou inexistente, além de não haver a preocupação de analisar as especificidades da linguagem literária, principalmente a questão da ambiguidade ; houve, também, muitas afirmações completamente erradas a respeito da linguagem literária ; no que diz respeito à forma do texto, foram observadas muitas incorreções gramaticais ( pontuação, concordância, acentuação, regência e seleção vocabular ) e texto muito truncado por problemas na paragrafação e no uso de conectivos. 3.3.2.3 Análise da Questão Discursiva 4 do Componente de Conhecimento Específico A Tabela 3.15 contém as informações relativas à questão discursiva 4 do conjunto de questões do componente de Conhecimento Específico.
O desempenho dos estudantes de todo o Brasil nesta questão foi inferior ao desempenho na questão de número 3.
A média geral do Brasil foi 23,1, sendo a menor média registrada na região Norte ( 15,6 ) e a maior na região Sudeste ( 27,0 ).
A nota máxima ( 100,0 ) foi atingida em todas as regiões.
A mediana foi 0,0 para o Brasil como um todo e para as regiões Norte, Nordeste e Centro-Oeste.
Na região Sudeste a mediana foi 20,0 e na Sul foi 15,0.
A nota mínima ( 0,0 ) foi a mesma em todas as regiões do Brasil.
O Gráfico 3.21, representa a distribuição de notas na questão discursiva 4, no componente de Conhecimento Específico para alunos do Bacharelado e da Licenciatura.
A maior frequência, tanto para estudantes de Bacharelado quanto para Licenciatura, ocorre na categoria “ em branco ”.
Observa -se que uma maior proporção de estudantes deixou essa questão em branco, comparando -se à questão discursiva 3.
Nota -se ainda um máximo local no intervalo ( 40 ; 50 ] para ambas as habilitações.
Os coeficientes de assimetria são positivos : 0,97 para a distribuição de notas dos concluintes de Bacharelado e 0,83 para os de Licenciatura.
Este comportamento se repete para ambos os cursos em todas as regiões, a menos da distribuição de notas do Bacharelado na região Centro-Oeste que tem coeficiente de assimetria negativo.
A assimetria positiva, longa cauda à direita do gráfico, é mais acentuada, com coeficiente de assimetria maior do que 1, para o Bacharelado da região Sul e para as duas habilitações da na região Norte. 3.3.2.4 Comentários sobre a correção das respostas à Questão Discursiva 4 A questão 4 apresenta temática e grau de dificuldade compatíveis com o nível de conhecimento esperado de um graduado em Letras.
Apesar disso, a análise das respostas evidenciou que apenas alguns estudantes demonstraram dominar de maneira satisfatória o conhecimento relativo a questões voltadas para o contexto sócio-histórico-cultural em que se inseriu o Movimento Modernista no Brasil, como também características de obras ligadas ao Modernismo.
A maioria absoluta abordou a questão apenas sob o ponto de vista temático, mas sem estabelecer qualquer relação com o movimento modernista.
Descreveram o subúrbio, atribuindo -lhe as características de “ violento ”, “ sombrio ”.
O aspecto social das favelas e subúrbios foi abordado por muitos estudantes, considerando as características dos seus moradores, sem estabelecer relação com o poema e a tela.
As respostas completas foram raras.
Em geral, a tendência dos estudantes foi realizar a descrição do poema e/ou da imagem.
Quando muito, reconheceram o movimento literário e passaram a discorrer sobre o Modernismo, sem vinculação com as obras.
Muitos apresentaram tendência a privilegiar o conteúdo, a temática do Modernismo, e, mais raramente, houve exploração dos dados formais das obras relacionados ao movimento artístico.
A forte incidência de respostas que evidenciaram grande dificuldade de leitura/análise do poema apresentado constituiu um preocupante indicador da falta de hábito de leitura de poemas por parte dos formandos de Letras.
Talvez isso decorra da dificuldade de leitura e interpretação de textos, recurso necessário para a análise formal das obras.
Muitas respostas estavam abaixo de qualquer nível esperado para concluintes de Letras.
Muitos demonstraram não saber o que é o Modernismo, não conhecer os poemas de Drummond e desconhecer a pintura de Tarsila.
Levando em conta que o Modernismo é trabalhado em sala de aula no Ensino Médio e nas aulas de Literatura nos cursos de graduação, é inconcebível que se apresentem respostas tão desconexas e, por vezes, absurdas.
O desconhecimento pleno, por exemplo, do que é o Modernismo os levou a vincular tanto o poema de Drummond quanto a tela de Tarsila ao Realismo, ao Naturalismo, ao Simbolismo, ao Iluminismo, ao Pós-Modernismo ou até ao Renascimento.
A dispersão de respostas se deveu, também, à inserção de um texto motivador que não tinha relação direta com o tema solicitado, o que levou um grande número de graduandos a dissertar sobre a diferença entre a obra literária e as artes plásticas, quanto à questão da construção da imagem ou da presença da imagem pronta.
Além da falta de domínio do conteúdo da questão, destacam -se as dificuldades de redação de muitos participantes, que produziram respostas, em geral, desconexas, com descrições rasas e com um sem número de dificuldades no que diz respeito ao uso da norma-padrão – problemas de ortografia, de concordância, de regência, de seleção vocabular, de pontuação, de utilização de mecanismos de coesão...
Evidencia -se, assim, em muitos textos, a falta de familiaridade com a norma-padrão, devido à abundância de desvios.
Parece que, para significativa parcela dos estudantes, há essa lacuna em sua formação.
Conclui -se, portanto, que pouquíssimos conseguiram de fato desenvolver os tópicos solicitados e relacioná -los produtivamente.
Houve um grande número de respostas que caíram no lugar-comum da descrição do “ subúrbio ” ; muitos casos de interpretação insatisfatória de um ou dos dois tópicos, havendo muitas respostas de pouquíssimas linhas.
De fato, poucos conseguiram desenvolver esses tópicos a contento, e menos ainda relacioná -los.
Destaca -se, também, o grande volume de respostas em branco.
Por outro lado, em respostas que apresentaram bom desempenho, tanto a forma quanto o conteúdo foram satisfatoriamente trabalhados : no que diz respeito ao conteúdo, identificaram -se os aspectos temáticos e formais característicos do Modernismo, comparando os dois textos ; quanto à forma do texto, não se foram observadas incorreções gramaticais nem ortográficas, houve paragrafação pertinente e adequado uso de conectivos.
Em respostas de desempenho médio, houve pequenas incorreções na forma ou no conteúdo : no que diz respeito ao conteúdo, percebeu -se uma tendência muito frequente de indicação da temática do Modernismo, sem, no entanto, exploração das questões formais das obras em questão, ou seja, os alunos não conseguiram relacionar os aspectos temáticos e formais ; no que diz respeito à forma do texto, observou -se incorreções gramaticais de pouca importância ( pontuação, concordância, acentuação e regência ) ou problemas na paragrafação e no uso de conectivos.
Já nas respostas de desempenho fraco, houve incorreções na forma ou no conteúdo : no que diz respeito ao conteúdo, os estudantes abordaram com bastante superficialidade o tema, ora detendo -se no texto motivador, ora fazendo afirmações genéricas sobre as questões sociais que envolvem subúrbio e favela no Brasil, muitas vezes confundindo os dois tipos de configuração urbana ; no que diz respeito à forma do texto, observaram -se muitas incorreções gramaticais ( pontuação, concordância, acentuação e regência ) e texto truncado por problemas na paragrafação e no uso de conectivos.
O aspecto social das favelas e subúrbios foi abordado por muitos concluintes, considerando as características dos seus moradores, sem estabelecer relação com o poema e a tela.
Houve respostas em que foram detectados equívocos conceituais, ao vincular as obras a outros movimentos literários, como o Naturalismo, ou atribuir a Machado de Assis a autoria do poema.
Além disso, apresentam inúmeros problemas de uso da norma-padrão.
Percebe -se que o estudante leu, interpretou, conhece o movimento literário, mas a apresentação do texto da resposta evidencia limitação não esperada do graduado em Letras : erros de concordância, de coerência, de coesão, de pontuação e de ortografia.
Grande parte dos concluintes abordou o tema Modernismo com superficialidade, revelando desconhecimento dos aspectos temáticos e formais da concepção filosófica da proposta Modernista.
Isso pôde ser notado pelas respostas que, muitas vezes, discorreram sobre características gerais do Modernismo.
Finalmente, em respostas que apresentaram desempenho insuficiente houve muitas incorreções na forma ou no conteúdo : no que diz respeito ao conteúdo, a abordagem dos textos foi incongruente, superficial ou tangencial ; no que diz respeito à forma do texto, observaram -se muitas incorreções gramaticais ( pontuação, concordância, acentuação, 3.3.2.5 Análise da Questão Discursiva 5 do Componente de Conhecimento Específico A Tabela 3.16 contém as informações relativas à questão discursiva 5 do conjunto do componente de Conhecimento Específico.
O desempenho dos estudantes nessa questão foi inferior ao das questões 3 e 4.
A nota média dos estudantes de todo o Brasil foi 19,8.
A maior média foi registrada na região Sudeste ( 23,4 ), enquanto a menor média foi registrada na região Norte ( 13,1 ).
Quanto à variabilidade das notas, o desvio padrão dos alunos do Brasil, como um todo, foi 28,2.
Enquanto o maior desvio foi encontrado na região Sudeste ( 30,1 ), o menor foi encontrado na região Norte ( 23,3 ).
A nota máxima ( 100,0 ) foi alcançada em todas as Grandes Regiões do Brasil.
A mediana e a nota mínima foram zero, também para todas as regiões, o que indica que pelo menos metade dos alunos de cada região tirou nota zero ou deixou a questão em branco.
O Gráfico 3.22 representa a distribuição das notas da questão discursiva 5 no componente de Conhecimento Específico para alunos do Bacharelado e da Licenciatura.
Mais uma vez, destaca -se o grande número de estudantes que deixaram a questão em branco, para as duas modalidades de curso.
Como nas questões discursivas do Componente de Conhecimento Específico anteriores, esta tem assimetria positiva.
Neste caso a assimetria é mais acentuada já que os coeficientes são maiores do que 1, o que explica a longa cauda à direita do gráfico.
Para alunos do Bacharelado o coeficiente de assimetria é de 1,07 e para Licenciatura de 1,14.
Este tipo de assimetria também ocorre em todas as regiões, com coeficientes maiores do que 1, para ambos os cursos, na maioria 3.3.2.5 Análise da Questão Discursiva 5 do Componente de Conhecimento Específico A Tabela 3.16 contém as informações relativas à questão discursiva 5 do conjunto do componente de Conhecimento Específico.
O desempenho dos estudantes nessa questão foi inferior ao das questões 3 e 4.
A nota média dos estudantes de todo o Brasil foi 19,8.
A maior média foi registrada na região Sudeste ( 23,4 ), enquanto a menor média foi registrada na região Norte ( 13,1 ).
Quanto à variabilidade das notas, o desvio padrão dos alunos do Brasil, como um todo, foi 28,2.
Enquanto o maior desvio foi encontrado na região Sudeste ( 30,1 ), o menor foi encontrado na região Norte ( 23,3 ).
A nota máxima ( 100,0 ) foi alcançada em todas as Grandes Regiões do Brasil.
A mediana e a nota mínima foram zero, também para todas as regiões, o que indica que pelo menos metade dos alunos de cada região tirou nota zero ou deixou a questão em branco.
O Gráfico 3.22 representa a distribuição das notas da questão discursiva 5 no componente de Conhecimento Específico para alunos do Bacharelado e da Licenciatura.
Mais uma vez, destaca -se o grande número de estudantes que deixaram a questão em branco, para as duas modalidades de curso.
Como nas questões discursivas do Componente de Conhecimento Específico anteriores, esta tem assimetria positiva.
Neste caso a assimetria é mais acentuada já que os coeficientes são maiores do que 1, o que explica a longa cauda à direita do gráfico.
Para alunos do Bacharelado o coeficiente de assimetria é de 1,07 e para Licenciatura de 1,14.
Este tipo de assimetria também ocorre em todas as regiões, com coeficientes maiores do que 1, para ambos os cursos, na maioria deles.
Coeficientes de assimetria de mesma ordem de \textbf{grandeza}, neste caso, sugerem que as formas dos histogramas devem ser semelhantes. 3.3.2.6 Comentários sobre a correção das respostas à Questão Discursiva 5 Trata -se de uma questão de pragmática.
Abordam -se as máximas conversacionais elaboradas pelo filósofo da linguagem Paul Grice.
Apesar das críticas em torno do fato de essas máximas governarem “ toda ” a comunicação [ p.ex., Jacob Mey aponta que elas foram pensadas no contexto de comunicação anglo-europeu ], a inserção dessa questão na prova do ENADE 2011 revela a expectativa de que os estudantes egressos de um curso de Letras conheçam as máximas e saibam analisar cenas da conversação a partir delas.
Através daquilo que formulou como “ Princípio de Cooperação ”, Grice procura explicar os critérios que orientam os falantes, em relação não apenas ao que se explicita na comunicação, mas principalmente àquilo que está implícito.
Trata -se de uma reflexão bastante complexa, em que se mostra que, apesar de haver regras que guiam a produção e interpretação de enunciados, dependendo da situação essas regras podem se apoiar ou, ao contrário, se contrariar.
O respeito a essas regras seria, portanto, instável, variável, sujeito a flutuação.
Com efeito, num momento, ele enuncia as quatro principais máximas ( qualidade, quantidade, modo e relevância – embora manifestando ao mesmo tempo consciência de que outras tantas máximas estariam sendo deixadas de lado, como a máxima da polidez : - Seja educado! ), porém logo adiante ele mostra que essas mesmas máximas podem entrar em concorrência.
Isso quer dizer que, em dadas situações, atender a uma máxima obriga o falante a transgredir outra máxima.
A questão solicitava que o estudante redigisse um texto dissertativo sobre duas máximas conversacionais de Grice, selecionadas no texto motivador.
Solicitava, também, que apresentasse exemplos de máximas que foram atendidas e máximas que foram violadas.
A questão valia 10,0 pontos.
O próprio enunciado já anunciava que exemplos de máximas atendidas valiam 5,0 pontos e exemplos de máximas violadas valiam 5,0 pontos, sem a previsão de pontuação para o trecho dissertativo solicitado na questão.
A questão 5 apresenta temática e grau de dificuldade compatíveis com o nível de conhecimento esperado de um graduando em Letras, embora esse assunto não conste, normalmente, da grade curricular dos cursos de graduação em Letras no Brasil, sendo abordado, preferencialmente, no âmbito dos programas de pós-graduação.
Assim sendo, o baixo rendimento dos estudantes poderia ser justificado, o que é comprovado por respostas como : “ não vi este conteúdo ” ; “ não lembro de ter visto essa matéria ”, “ quem é Greice ”?
Outra evidência de que esse conteúdo não é comumente abordado na graduação pode ser observada pela análise de respostas que tentam vincular as máximas ao ensino de qualidade ( em referência à máxima da qualidade ), à necessidade de se dizer a verdade e não mentir, pois “ muitas pessoas são falsas umas com as outras ” ( em referência à máxima da qualidade ) e às boas maneiras, ser educado ( em referência à máxima do modo ).
É preciso destacar, também, que houve confusão entre “ violação das máximas de Grice ” e “ violação da norma culta ou da norma-padrão ”.
Dessa forma, houve menções como : “ não se pode infringir a norma culta da língua, sob o risco de não sermos compreendidos ”.
Soma -se a isso a verificação de que há equívocos relacionados aos conceitos de norma culta e norma-padrão, bem como aos conhecimentos sobre variação linguística. interpretação de enunciados, dependendo da situação essas regras podem se apoiar ou, ao contrário, se contrariar.
O respeito a essas regras seria, portanto, instável, variável, sujeito a flutuação.
Com efeito, num momento, ele enuncia as quatro principais máximas ( qualidade, quantidade, modo e relevância – embora manifestando ao mesmo tempo consciência de que outras tantas máximas estariam sendo deixadas de lado, como a máxima da polidez : - Seja educado! ), porém logo adiante ele mostra que essas mesmas máximas podem entrar em concorrência.
Isso quer dizer que, em dadas situações, atender a uma máxima obriga o falante a transgredir outra máxima.
A questão solicitava que o estudante redigisse um texto dissertativo sobre duas máximas conversacionais de Grice, selecionadas no texto motivador.
Solicitava, também, que apresentasse exemplos de máximas que foram atendidas e máximas que foram violadas.
A questão valia 10,0 pontos.
O próprio enunciado já anunciava que exemplos de máximas atendidas valiam 5,0 pontos e exemplos de máximas violadas valiam 5,0 pontos, sem a previsão de pontuação para o trecho dissertativo solicitado na questão.
A questão 5 apresenta temática e grau de dificuldade compatíveis com o nível de conhecimento esperado de um graduando em Letras, embora esse assunto não conste, normalmente, da grade curricular dos cursos de graduação em Letras no Brasil, sendo abordado, preferencialmente, no âmbito dos programas de pós-graduação.
Assim sendo, o baixo rendimento dos estudantes poderia ser justificado, o que é comprovado por respostas como : “ não vi este conteúdo ” ; “ não lembro de ter visto essa matéria ”, “ quem é Greice ”?
Outra evidência de que esse conteúdo não é comumente abordado na graduação pode ser observada pela análise de respostas que tentam vincular as máximas ao ensino de qualidade ( em referência à máxima da qualidade ), à necessidade de se dizer a verdade e não mentir, pois “ muitas pessoas são falsas umas com as outras ” ( em referência à máxima da qualidade ) e às boas maneiras, ser educado ( em referência à máxima do modo ).
É preciso destacar, também, que houve confusão entre “ violação das máximas de Grice ” e “ violação da norma culta ou da norma-padrão ”.
Dessa forma, houve menções como : “ não se pode infringir a norma culta da língua, sob o risco de não sermos compreendidos ”.
Soma -se a isso a verificação de que há equívocos relacionados aos conceitos de norma culta e norma-padrão, bem como aos conhecimentos sobre variação linguística.
Entretanto, pode -se argumentar que, apesar desse assunto não ser tratado na maioria dos cursos de Letras ( realmente, a disciplina Pragmática consta em poucas grades curriculares do Brasil ), o enunciado da questão oferece informações suficientes para a compreensão do assunto – e qualquer usuário proficiente do português pode recorrer à sua intuição e responder ao item.
Desconsiderando -se o grande número de respostas deixadas em branco, houve um desempenho de fraco a médio.
A maioria dos estudantes transcreveu do enunciado a explicação sobre as máximas escolhidas.
Alguns deles atenderam ao todo da questão e outros, ao contrário, tangenciaram a resposta.
A julgar pela maioria das respostas, houve uma grande dificuldade de elaborar um texto de base dissertativa, revelando que os formandos em Letras não estão desenvolvendo estratégias de reflexão sobre assuntos relevantes a profissionais que terão a língua e a literatura como instrumentos de trabalho.
Nessas respostas, não se verificou a preocupação de discutir o problema, mas sim de fazer afirmações destituídas de argumentos que as comprovem.
Na verdade, o aluno apresentou opiniões prontas.
Houve poucas respostas que contemplaram adequadamente o conteúdo exigido : explicação das máximas conversacionais e exemplificação com diálogos em que as máximas fossem respeitadas e violadas.
Normalmente, os exemplos foram explicações de possível violação : “ um exemplo de máxima violada é o caso de alguém perguntar as horas e o outro responder uma coisa que não tem nada a ver ”.
Por estarem tendo contato pela primeira vez com esse conteúdo, alguns graduandos tentaram responder à questão com base apenas na leitura do texto teórico motivador.
Boa parte dos que escolheram tratar da máxima de qualidade, por exemplo, entendeu que o seguinte trecho do enunciado : “ a contribuição conversacional do falante deve ser a mais verdadeira possível ” significava que o falante só deveria dizer verdades incontestáveis, verificáveis empiricamente.
No entanto, significa apenas que, em uma situação de comunicação, o destinatário costuma partir do princípio de que o enunciador acredita naquilo que ele enuncia.
Outro ponto da questão 5 que favoreceu uma compreensão extremamente rasa do assunto foi a definição da máxima da relevância : “ a contribuição conversacional do falante deve ser pertinente em relação ao objeto da conversa ”.
Atendendo ao que se lhes pedia, vários estudantes tentaram fornecer exemplos de violação dessa máxima, que, no entanto, revelavam sua incompreensão da mesma.
São diálogos como o seguinte : - Você vai sair hoje à noite? - Ah, hoje eu estou muito cansado.
Em termos de encadeamento, existe uma perfeita relevância dessa resposta em relação à pergunta, relevância esta que, aliás, costuma estar presente em qualquer enunciado.
Costuma -se dizer, inclusive, que a relevância é uma constante ; o sentido do enunciado, este sim, é variável.
Sem a presença do princípio de “ relevância ”, como poderíamos interpretar a resposta fornecida no exemplo logo acima?
Essa mesma relevância orienta a interpretação de sequências como : - Estou ficando com sono. - Vou -lhe trazer um café. - Você sabe que horas são? - Desculpe meu atraso.
O pneu do carro furou.
Tive de vir a pé.
Trata -se de diálogos banais, perfeitamente coerentes e compreensíveis, que não emperram de forma alguma a máquina conversacional.
É, portanto, inegável que o enunciado da questão 5 levou boa parte dos estudantes a uma lamentável confusão ( entre o plano dos fatos “ reais ” e o plano discursivo ; entre implicitação e explicitação ), o que prejudicou a avaliação pretendida.
De modo geral, foi uma questão com \textbf{elevadíssima} taxa de respostas em branco, ou seja, os formandos sequer se deram ao trabalho de respondê -la.
E, quando havia resposta, em muitos casos eram desconsideradas, já que aproveitavam o espaço para dizer que não haviam tratado daquele conteúdo em sua graduação ou que o tempo era insuficiente para o tamanho da prova do ENADE como um todo ; ou eram zeradas, já que os graduandos, no esforço de atender ao solicitado, mesmo sem conhecimento, tentaram formular qualquer resposta – nesse sentido, recorriam a conteúdos ligados à sociolinguística, à análise da conversação e do discurso, às teorias do texto, entre outras.
Em relação ao desenvolvimento do conteúdo, portanto, pode -se concluir que a questão 5 não concorreu para a efetiva avaliação do desempenho dos cursos de graduação em Letras.
Um importante aspecto a destacar é o mau desempenho da grande maioria dos alunos em relação à estrutura formal do texto produzido ( mecanismos de coesão e utilização do padrão formal culto ), o que é extremamente preocupante ao se levar em conta que são formandos em Letras.
Foram muito comuns os casos de erros de concordância, pontuação, acentuação, regência, seleção vocabular e estruturação frasal, com uso inadequado ou ausência de conectivos entre parágrafos e entre frases.
Os textos demonstraram falta de domínio de mecanismos básicos de textualidade e do registro padrão da língua.
Nas respostas que apresentaram bom desempenho, tanto a forma quanto o conteúdo foram satisfatoriamente trabalhados : no que diz respeito ao conteúdo, o estudante desenvolveu uma pequena explicação sobre as duas máximas escolhidas, apresentou exemplos de seu atendimento e de sua violação, com transcrição de diálogos fictícios ou com a descrição de situações de conversação ; quanto à forma do texto, não foram observadas incorreções gramaticais nem ortográficas, houve paragrafação pertinente e adequado uso de conectivos.
Em respostas que apresentaram desempenho mediano, houve pequenas incorreções na forma ou no conteúdo : no que diz respeito ao conteúdo, os estudantes não atenderam a todos os itens solicitados no enunciado, abordando apenas uma das máximas, sem a explicitação dissertativa, ou apresentando apenas exemplos de atendimento ou de violação ; no que diz respeito à forma do texto, foram observadas incorreções gramaticais de pouca importância ( pontuação, concordância, acentuação e regência ) ou problemas na paragrafação e no uso de conectivos.
Em respostas desempenho fraco, houve incorreções na forma ou no conteúdo : no que diz respeito ao conteúdo, alguns estudantes dissertaram de um modo geral sobre as máximas, sem apresentar exemplos ; ou, ao contrário, exemplificaram as máximas, mas não explicaram o que significa cada máxima selecionada ; no que diz respeito à forma do texto, observou -se muitas incorreções gramaticais ( pontuação, concordância, acentuação e regência ) e texto truncado por problemas na paragrafação e no uso de conectivos.
Houve um domínio fraquíssimo dos mecanismos de estruturação textual e das regras do padrão formal da língua, como concordância e pontuação.
Quanto ao conteúdo, o estudante tentou dissertar sobre as máximas, mas não exemplificou.
Em desempenhos insuficientes, houve muitas incorreções na forma ou no conteúdo : no que diz respeito ao conteúdo, a abordagem das máximas conversacionais foi praticamente nula ou incompreensível, o que revela total desconhecimento do conteúdo teórico exigido na questão ; no que diz respeito à forma do texto, constatou -se incorreções gramaticais em pontuação, concordância, acentuação, regência e seleção vocabular.
Além disso, textos muito truncados por problemas na paragrafação e no uso de conectivos. 3.3.3 Considerações Finais da Banca sobre a correção das questões discursivas do Componente de Conhecimento Específico Ao final do processo de correção das três questões, pode -se afirmar que poucos estudantes conseguiram desenvolver com aprofundamento o tema solicitado e utilizar o padrão formal da língua.
Com base nos resultados obtidos, constatou -se que a grande maioria dos que estão se formando nos cursos de graduação em Letras no nosso país não domina os conteúdos cobrados nas três questões discursivas, conteúdos esses considerados básicos para esse perfil profissional.
As respostas analisadas não nos permitem afirmar com segurança como esses conteúdos foram trabalhados na graduação, a não ser pelo sucesso de parte dos concluintes, pelo insucesso de outra parte e por algumas manifestações do tipo “ este assunto não foi tratado na faculdade ”.
Um fator que prejudica consideravelmente a avaliação do desempenho dos estudantes é o grande número de provas em branco, o que leva a especular sobre as possíveis causas : ou não sabem responder à questão ou apenas se recusam a respondê -la, por motivos diversos, como a grande extensão da prova ( comentada por muitos ) ou a negação da validade do processo de avaliação desenvolvido pelo ENADE.
Pela análise das questões, percebe -se que os estudantes apresentam dificuldades tanto no plano do conteúdo quanto no plano da expressão.
Muitos não dominam regras básicas de estruturação textual e produzem textos desconexos, sem qualquer pensamento organizado.
O domínio da norma-padrão e a capacidade de argumentação, em termos de estilo e estrutura textual, mostrou -se muito aquém do desejado na maioria das respostas.
É digno de nota que os estudantes que demonstraram menor compreensão do conteúdo da questão foram os mesmos que apresentaram pior desempenho na produção textual.
Isso é um indicador relevante de que a capacidade de raciocinar por escrito é diretamente proporcional à capacidade de leitura.
Tudo isso leva a refletir sobre a necessidade de repensar o estudo da Língua Portuguesa e da Literatura no Ensino Médio e na Universidade.
É necessário, por exemplo, que os estudantes tenham contato com as obras literárias, por meio de leituras e análises de obras representativas da literatura em língua portuguesa, e não apenas conheçam as características dos movimentos literários.

% matched lemmas: creche, elevado, grandeza, isenção, negrito, vicioso
\end{textsample}
